\section{Conclusion}\label{sec:conclusion}

We have described a text-associated multi-morphology motion corpus and a
physics-based imitation system that preserves the correspondence between
source clips, generated references, and simulated bodies. The construction
combines AMASS/HumanML3D provenance, HUMOS generation, SMPLSim assets,
conservative reference refinement, and deterministic clip-grouped splits.
The controller uses temporal slot/type attention with morphology information
in both actor and critic and is trained directly across body configurations.

Offline evaluation of the selected checkpoint gives 98.47\% success on
validation and 99.14\% on the untouched test split for the one-shape-per-clip
trained-body protocol, and the eight-run development grid selects slot/type
attention without adaLN-Zero. Residual failures on both splits are dominated
by ground-contact transitions and whole-body acrobatics. These results do not
establish all-shape coverage or unseen-body generalization. Two studies are
specified but not run: an unseen-morphology evaluation over body shapes
inside the trained coefficient range, and an analysis of how morphology and
motion type jointly set physical control demand -- whether dynamic motions
are more shape-sensitive than quasi-static ones.
